	\section{Introducción}
	El presente documento reúne el avance del Trabajo Terminal 2026-B009, desarrollado durante el
	semestre 2025-2 en colaboración con la Defensoría de los Derechos Politécnicos (DDP) del Instituto
	Politécnico Nacional (IPN), con el objetivo de identificar áreas de oportunidad en sus procesos
	internos y proponer una solución tecnológica que los optimice.
	
	El Instituto Politécnico Nacional es una institución educativa pública mexicana consolidada como una
	de las principales casas de estudio del país, la cual ofrece una amplia gama de programas académicos
	en diversas áreas del conocimiento \cite{ipn_ley_organica}.
	
	La transformación digital ha modificado profundamente la manera en que las instituciones administran
	sus procesos internos, resguardan información sensible y prestan servicios a sus usuarios. En este
	contexto, las organizaciones educativas enfrentan retos nuevos relacionados con la atención de
	quejas, la protección de los derechos de su comunidad y la necesidad de contar con mecanismos
	tecnológicos capaces de garantizar seguridad, trazabilidad y eficiencia operativa. Dentro del IPN,
	la Defensoría de los Derechos Politécnicos desempeña un papel fundamental como órgano encargado de
	promover, proteger y defender los derechos de estudiantes, docentes, personal administrativo y demás
	integrantes de la comunidad politécnica \cite{ipn_acuerdo622, ipn_manual_organizacion_ddp}.
	
	Actualmente, buena parte de los procedimientos de recepción, análisis y seguimiento de quejas
	depende de procesos manuales o parcialmente digitalizados, lo que ocasiona diversos problemas
	operativos \cite{ipn_manual_procedimientos_ddp}. Entre las limitaciones identificadas destacan la
	saturación administrativa, la dependencia del criterio humano en la etapa de primer contacto, la
	dificultad para dar seguimiento oportuno a los expedientes y la ausencia de mecanismos inteligentes
	que automaticen la clasificación y canalización de casos. Estas problemáticas provocan retrasos en
	la atención, inconsistencias en la interpretación normativa y una capacidad institucional limitada
	para responder al incremento constante de solicitudes.
	
	A partir de este escenario surge la necesidad de desarrollar una plataforma web especializada en la
	gestión digital de quejas y denuncias dentro de la DDP-IPN. La propuesta que aquí se presenta busca
	diseñar e implementar una solución tecnológica moderna que optimice el flujo operativo institucional
	mediante la integración de herramientas de procesamiento de lenguaje natural, aprendizaje
	automático, arquitectura de microservicios y mecanismos avanzados de seguridad informática. El
	propósito no se limita a digitalizar formularios o expedientes, sino a construir un ecosistema
	integral capaz de asistir al personal jurídico y administrativo en la toma de decisiones sobre la
	admisión y clasificación de quejas.
	
	El proyecto se fundamenta en un análisis previo del estado del arte de plataformas gubernamentales,
	académicas e internacionales orientadas a la gestión de denuncias digitales
	\cite{sfp_sidec,cdmx_denuncia_digital,whistlelink,eqs_integrity_line,navex_whistleb,unam_ddu_denuncia}.
	Dicho análisis permitió identificar que las soluciones existentes presentan limitaciones importantes en
	escalabilidad, automatización inteligente, protección de identidad, interoperabilidad institucional
	y especialización universitaria. Mientras algunas priorizan el cumplimiento normativo y la
	formalización administrativa, otras se enfocan en la anonimización o la gestión corporativa, y dejan
	vacíos relevantes en contextos educativos donde convergen necesidades jurídicas, administrativas y
	de protección de derechos humanos. De ahí que se identificara una oportunidad tecnológica para
	desarrollar una solución adaptada a las necesidades operativas de la DDP.
	
	Uno de los principales aportes de este trabajo consiste en la incorporación de un clasificador
	inteligente de quejas basado en técnicas de Procesamiento de Lenguaje Natural (PLN) y aprendizaje
	supervisado \cite{cortes1995svm, breiman2001rf, reimers2019sbert}. Este componente asiste al
	personal de primer contacto en la determinación de competencia o incompetencia de una queja, y toma
	como referencia el marco normativo establecido en el Acuerdo Número Extraordinario 622 de la Gaceta
	Politécnica \cite{ipn_acuerdo622}. Hoy en día esa decisión depende del análisis manual que realiza
	el personal administrativo y jurídico, lo cual genera variabilidad en los criterios de
	interpretación y un consumo considerable de tiempo operativo.
	
	Además del componente de inteligencia artificial, el proyecto contempla una arquitectura moderna
	basada en microservicios, orientada a garantizar modularidad, escalabilidad y mantenibilidad del
	sistema \cite{aws_microservicios, villamizar_microservices}. Entre los módulos principales destacan
	el portal del quejoso, el módulo administrativo, el expediente digital, el sistema de seguridad y
	control de acceso, y el clasificador inteligente. Esta arquitectura desacopla las funcionalidades
	críticas y facilita futuras ampliaciones sin afectar la operación general de la plataforma.
	
	Como parte del trabajo realizado se mantuvo una interacción constante con la DDP mediante reuniones,
	levantamiento de requerimientos y validación de procesos operativos. Ello permitió modelar el flujo
	real de atención de quejas con diagramas BPMN, identificar reglas de negocio críticas y diseñar
	soluciones alineadas con las necesidades institucionales. También se elaboraron prototipos
	funcionales, mockups, modelos de base de datos, la implementación de backend y frontend del
	microservicio de atención al quejoso, así como pruebas funcionales y experimentales para validar el
	comportamiento del sistema.
	
	En conjunto, este trabajo presenta una propuesta integral para modernizar la gestión de quejas y
	denuncias dentro de la Defensoría de los Derechos Politécnicos del IPN. La integración de
	inteligencia artificial, arquitecturas escalables y mecanismos de seguridad avanzada sienta las
	bases de una plataforma institucional capaz de optimizar procesos administrativos, reducir tiempos
	de respuesta y fortalecer la protección de los derechos de la comunidad politécnica. En los
	capítulos siguientes se desarrollan los fundamentos teóricos, metodológicos y técnicos que sustentan
	el diseño, la implementación y la validación de la solución propuesta.
